\documentclass[11pt,a4paper]{article}

%================================================
% ENCODAGE ET LANGUE
%================================================

\usepackage[utf8]{inputenc}
\usepackage[T1]{fontenc}
\usepackage[french]{babel}

%================================================
% MISE EN PAGE
%================================================

\usepackage{geometry}
\geometry{margin=1.55cm}

%================================================
% PACKAGE GTOLI
%================================================

\usepackage{../styles/gtoli}

%================================================
% INFORMATIONS DU DOCUMENT
%================================================

\title{\textbf{Cours -- 6e secondaire}\\[0.2cm]
\large Titre du chapitre}

\author{GToli}
\date{\today}

%================================================
% DOCUMENT
%================================================

\begin{document}

\maketitle

%================================================
% OBJECTIFS
%================================================

\begin{cadreobjectif}
À la fin de ce chapitre, l’élève doit être capable de :
\begin{itemize}
    \item mobiliser les notions avec autonomie ;
    \item construire un raisonnement structuré ;
    \item justifier les choix méthodologiques ;
    \item résoudre des problèmes complexes ;
    \item articuler plusieurs notions dans une même démarche ;
    \item rédiger une solution rigoureuse et synthétique ;
    \item préparer les exigences de l’enseignement supérieur.
\end{itemize}
\end{cadreobjectif}

%================================================
\section{Enjeux du chapitre}
%================================================

\begin{cadrebleu}{Mise en perspective}
Ce chapitre s’inscrit dans une logique de consolidation et de préparation à l’enseignement supérieur.

Il vise à développer :
\begin{itemize}
    \item l’autonomie dans le choix des méthodes ;
    \item la rigueur dans la rédaction ;
    \item la capacité à relier plusieurs notions ;
    \item la lecture critique des résultats ;
    \item la maîtrise des outils conceptuels et techniques.
\end{itemize}
\end{cadrebleu}

%================================================
\section{Prérequis et connaissances mobilisables}
%================================================

\begin{cadregris}{Prérequis}
Avant d’aborder ce chapitre, l’élève doit être capable de :
\begin{itemize}
    \item reconnaître les notions fondamentales associées ;
    \item appliquer les propriétés de base ;
    \item manipuler les notations avec précision ;
    \item utiliser les méthodes vues antérieurement ;
    \item rédiger une justification courte mais correcte.
\end{itemize}
\end{cadregris}

%================================================
\section{Cadre théorique}
%================================================

\begin{cadretheorie}
Cette section présente les éléments conceptuels indispensables :

\begin{itemize}
    \item définitions formelles ;
    \item propriétés ;
    \item théorèmes ;
    \item conditions d’application ;
    \item conventions de notation ;
    \item liens avec d’autres chapitres ;
    \item limites de validité des méthodes.
\end{itemize}

\medskip

\textbf{Exigence de 6e :} la théorie doit pouvoir être mobilisée dans des situations non immédiatement reconnaissables.
\end{cadretheorie}

%================================================
\section{Méthode experte}
%================================================

\begin{cadremethode}
Pour traiter un problème de niveau 6e secondaire :

\begin{enumerate}
    \item lire la question et identifier la nature du problème ;
    \item repérer les notions potentiellement mobilisables ;
    \item déterminer les contraintes et conditions d’application ;
    \item choisir une stratégie de résolution ;
    \item organiser le raisonnement ;
    \item effectuer les calculs ou démonstrations nécessaires ;
    \item contrôler la cohérence mathématique ou scientifique ;
    \item interpréter le résultat dans le contexte ;
    \item rédiger une conclusion précise.
\end{enumerate}
\end{cadremethode}

%================================================
\section{Carte des méthodes}
%================================================

\begin{cadregris}{Choisir la bonne méthode}
\renewcommand{\arraystretch}{1.35}
\begin{tabularx}{\textwidth}{>{\bfseries}p{3.5cm}X X}
\toprule
Situation rencontrée & Question à se poser & Méthode possible \\
\midrule
Cas simple & La méthode directe suffit-elle ? & Application immédiate \\
Cas avec condition & Y a-t-il une restriction ou une hypothèse ? & Vérifier les conditions \\
Cas complexe & Plusieurs notions sont-elles liées ? & Décomposer le problème \\
Cas de transfert & La situation ressemble-t-elle à un cas déjà vu ? & Adapter la méthode \\
Cas de synthèse & Faut-il articuler plusieurs chapitres ? & Construire une stratégie globale \\
\bottomrule
\end{tabularx}
\end{cadregris}

%================================================
\section{Exemple expert commenté}
%================================================

\begin{cadreexemple}
\textbf{Énoncé :}

Insérer ici un exercice représentatif du niveau attendu en 6e secondaire.

\medskip

\textbf{Analyse stratégique :}
\begin{itemize}
    \item Quelle est la nature du problème ?
    \item Quelles notions peuvent intervenir ?
    \item Quelle stratégie semble la plus efficace ?
    \item Quelles erreurs faut-il anticiper ?
\end{itemize}

\medskip

\textbf{Résolution :}

La résolution doit montrer non seulement les calculs, mais aussi les choix de méthode.
\end{cadreexemple}

%================================================
\section{Rédaction attendue}
%================================================

\begin{cadrevert}{Rigueur de rédaction}
En 6e secondaire, une solution complète doit être :
\begin{itemize}
    \item structurée ;
    \item justifiée ;
    \item concise ;
    \item mathématiquement ou scientifiquement correcte ;
    \item interprétée lorsque le contexte l’exige.
\end{itemize}

\medskip

L’objectif est de préparer l’élève à une rédaction compatible avec les exigences de l’enseignement supérieur.
\end{cadrevert}

%================================================
\section{Points de vigilance}
%================================================

\begin{cadreattention}
Les difficultés fréquentes sont :
\begin{itemize}
    \item appliquer une méthode sans vérifier les conditions ;
    \item confondre résultat numérique et interprétation ;
    \item rédiger uniquement des calculs sans justification ;
    \item ne pas identifier la stratégie globale ;
    \item perdre de vue la question posée ;
    \item oublier de vérifier la cohérence finale.
\end{itemize}
\end{cadreattention}

%================================================
\section{Exercice d’application autonome}
%================================================

\begin{cadreenonce}
\textbf{Exercice 1 :}

Insérer ici un exercice où l’élève doit choisir lui-même la méthode pertinente.
\end{cadreenonce}

\begin{cadreaide}
Questions stratégiques :
\begin{itemize}
    \item Quelle est la notion principale ?
    \item Quelle méthode paraît adaptée ?
    \item Quelles conditions doivent être vérifiées ?
    \item Comment vérifier la cohérence du résultat ?
\end{itemize}
\end{cadreaide}

%================================================
\section{Solution structurée}
%================================================

\begin{cadresolution}
La solution attendue doit comprendre :
\begin{enumerate}
    \item analyse du problème ;
    \item choix de la stratégie ;
    \item développement ;
    \item justification ;
    \item vérification ;
    \item conclusion.
\end{enumerate}
\end{cadresolution}

%================================================
\section{Exercices de transfert et de synthèse}
%================================================

\begin{cadreactivite}
\begin{enumerate}
    \item Exercice d’application autonome.
    \item Exercice de transfert.
    \item Exercice combinant plusieurs notions.
    \item Exercice de synthèse.
    \item Exercice de préparation à l’enseignement supérieur.
\end{enumerate}
\end{cadreactivite}

%================================================
\section{Mini-problème intégrateur}
%================================================

\begin{cadreenonce}
\textbf{Problème intégrateur :}

Insérer ici une situation complexe nécessitant plusieurs étapes de raisonnement.

\medskip

L’objectif est d’obliger l’élève à construire une stratégie complète et non à appliquer mécaniquement une formule.
\end{cadreenonce}

%================================================
\section{Synthèse conceptuelle}
%================================================

\begin{cadresynthese}
À retenir :

\begin{itemize}
    \item notion centrale : \dotfill
    \item méthodes disponibles : \dotfill
    \item conditions d’application : \dotfill
    \item liens avec d’autres chapitres : \dotfill
    \item stratégie de résolution : \dotfill
    \item erreur critique à éviter : \dotfill
\end{itemize}
\end{cadresynthese}

%================================================
\section{Auto-évaluation avancée}
%================================================

\begin{cadregris}{Niveau de maîtrise}
\begin{itemize}
    \item[$\square$] Je reconnais la nature du problème.
    \item[$\square$] Je choisis une méthode adaptée.
    \item[$\square$] Je vérifie les conditions d’application.
    \item[$\square$] Je justifie mon raisonnement.
    \item[$\square$] Je peux résoudre un exercice de transfert.
    \item[$\square$] Je peux articuler plusieurs notions.
    \item[$\square$] Je peux rédiger une solution concise et rigoureuse.
\end{itemize}
\end{cadregris}

%================================================
\section{Préparation à l’enseignement supérieur}
%================================================

\begin{cadreorange}{Conseils méthodologiques}
Pour se préparer efficacement :
\begin{enumerate}
    \item refaire les exercices sans regarder les solutions ;
    \item expliquer oralement la méthode utilisée ;
    \item comparer plusieurs stratégies possibles ;
    \item rédiger des solutions complètes ;
    \item identifier les erreurs récurrentes ;
    \item travailler la précision du vocabulaire ;
    \item apprendre à justifier, et pas seulement à calculer.
\end{enumerate}
\end{cadreorange}

\end{document}